\scanpage{227}
\margin{en vertu de la conjecture de Hodge}
D'autre part, \uline{quels} sont les foncteurs multiplicatifs de Motifs, à valeurs dans les structures de Hodge, \inserted{(\unsure{inclure} les motifs sur $\mathbf C$!)} qui sont associés à des places archimédiennes ? Plus généralement, dans quelle mesure un morphisme de préschémas \inserted{$X\mathrel{\text{---}}Y$} est connu par le foncteur sur les motifs \inserted{(loc. cit)}, et dans quelle mesure un foncteur \inserted{multiplicatif} peut-il être pris arbitrairement ?!?

C'est lié au Pb de \uline{construire des schémas de modules} pour les motifs polarisés; si sur $Y$ \inserted{\unsure{se donne} un motif} polarisé \unsure{global} tel que $Y\unsure{\longrightarrow|\mathrel{-\!\circ}}\text{Modules}$ est une immersion \inserted{\unsure{voire (resp. une)}}) \deleted{\unsure{est}} alors est \unsure{anabélien}. Kif kif si $Y$ est \uline{\unsure{l'un|lisse}} des préschémas ayant cette belle propriété. \unsure{Marche} en particulier si $Y$ est le spectre d'un corps\ldots\ldots
\mspage{226}
